% Source presentations for register row N-c-008.

\subsection{The monadic presentation of B\"orger--Gr\"adel--Gurevich}
\label{sec:p09-presentation-n-c-008}

In the source's bracket notation, the relational monadic class is
\([\mathrm{all},(\omega)]\): quantifier prefixes are unrestricted, every
nonlogical predicate symbol is unary, and there are no function symbols.  A
superscript \({}^{=}\) admits equality.  The book separately calls the language
which may also contain unary function symbols the \emph{full monadic class}
\cite[pp.~249--250]{BorgerGradelGurevich1997}.  Thus ``monadic'' does not by
itself decide whether unary functions or equality are present.  The handbook's
lookup row fixes the narrower equality-free, function-free convention; the
source presentation here records the adjacent variants explicitly.

Let \(\psi\) be a satisfiable relational monadic sentence of positive
quantifier rank \(q\), using unary predicates \(P_1,\ldots,P_m\), with equality
allowed.  (At quantifier rank (0), a satisfiable sentence in this
constant-free positive-arity vocabulary has the trivial one-element case.)  The
\emph{colour} of an element \(a\) is the Boolean vector
\[
 \bigl(\mathbf 1_{P_1(a)},\ldots,\mathbf 1_{P_m(a)}\bigr)
 \in\{0,1\}^m.
\]
This is an eliminable abbreviation for the element's complete unary atomic
type; at most \(2^m\) colours occur.  From any model of \(\psi\), retain at
most \(q\) elements of each colour and take the induced substructure.  In the
\(q\)-move Ehrenfeucht--Fra\"iss\'e game, Duplicator always matches a previously
unused element of the same colour: if fewer than \(q\) elements of that colour
exist, all were retained; if at least \(q\) exist, fewer than \(q\) moves have
been played.  Equality and every unary predicate are preserved.  Therefore the
reduced structure still satisfies \(\psi\) and has size at most
\[
 q\,2^m
\]
\cite[prop.~6.2.1]{BorgerGradelGurevich1997}.

This complete small-model proof gives decidability of relational monadic
satisfiability.  Without equality, Exercise~6.2.3 asks the reader to sharpen
the bound to \(2^m\); the source does not present that exercise as a completed
proof.  For equality-bearing formulas of length \(n\), Corollary~6.2.4 derives
a minimal-model bound \(2^{O(n/\log n)}\), followed by the corresponding
nondeterministic satisfaction procedure
\cite[pp.~250--251]{BorgerGradelGurevich1997}.  Exercise~6.2.5 likewise asks for
the extension to second-order formulas in which every free and bound predicate
is monadic.  That Behmann--Skolem extension remains exercise-level in this
presentation.

Unary functions require a different construction.  Proposition~6.2.7 first
replaces each atom \(P(f(t))\) by a fresh unary predicate together with an
equivalence \(P(f(x))\leftrightarrow Q(x)\), then replaces the remaining unary
function requirements by existential witnesses.  The result is equisatisfiable
with a relational monadic formula over the same domains and yields a model of
size \(2^{O(n)}\) for an input of length \(n\).  Proposition~6.2.15 shows that
this larger scale is substantive: a length-\(O(n)\) formula with one unary
function can force at least \((n+1)2^n\) elements by a chain feeding a long
counter cycle \cite[pp.~251--257]{BorgerGradelGurevich1997}.  The function
elimination is therefore not the same proof as retaining representatives of
unary colours.

\paragraph{L\"owenheim and Behmann source boundaries.}
The 1997 book's historical notes report that L\"owenheim formulates the 1915
result through validity, treats equality without functions, recognizes
infinity axioms, and reduces general validity to formulas with binary
predicates.  They report that Behmann's 1922 elimination method starts with the
equality-free monadic case and extends to equality and monadic second order
\cite[pp.~310--311]{BorgerGradelGurevich1997}.  Those are later historical
reports, not transcriptions of the originals.  A complete scan of Behmann's
paper establishes its six-part, seventeen-section elimination architecture,
but without reliable text it does not establish the exact symbols or proof
steps.  A readable L\"owenheim source was likewise unavailable.  The modern
colour/EF proof and unary-function construction above are therefore attributed
only to the 1997 presentation, while the two original presentations remain
open rather than being filled by translation without source pages
\cite{Lowenheim1915,Behmann1922}.
